\section{The different and the geometric genus}\label{sec:ramification}

We continue to use geometric constants for all local calculations.
The zero-place index was determined without the group induction in
Lemma~\ref{lem:zero}. The extension there is tame, so
\[
 d_{0,n}=e_{0,n}-1=2^n-1.
\]
At infinity, it is simpler to calculate the different first in
one root field and then pass to the Galois completion.

\begin{proposition}[Infinity different]\label{prop:different}
The different exponent of $\Lbar_n/\kk(t)$ at every place
above infinity is $3^{n+1}-2$.
\end{proposition}

\begin{proof}
Choose a root $X$ with $t=f^{\circ n}(X)$. Its function field
is the rational field $\kk(X)$. With $u=1/X$, the local
reciprocal map at infinity is
\[
 h(u)=\frac1{f(1/u)}=\frac{u^3}{1+au},
 \qquad h'(u)=\frac{-au^3}{(1+au)^2}.
\]
Thus $\ord_u h=3$ and $\ord_u h'=3$.
The base uniformizer is $t^{-1}=h^{\circ n}(u)$.
The root completion $\kk((u))/\kk((t^{-1}))$ is separable,
totally ramified of degree $3^n$. The derivative chain rule gives
\begin{align}
 \ord_u\bigl((h^{\circ n})'(u)\bigr)
   &=\sum_{j=0}^{n-1}\ord_u h'\bigl(h^{\circ j}(u)\bigr)\notag\\
   &=3\sum_{j=0}^{n-1}3^j
     =\frac{3(3^n-1)}2.\label{eq:root-different}
\end{align}
By the local different formula, this is the different exponent
of the root completion over the base completion.

Choose a place of $\Lbar_n$ extending this root-field place.
Its ramification index over the base is $2\cdot3^n$, already
proved in Section~\ref{sec:as}. All residue fields are $\kk$,
so the Galois completion has degree two over the root completion.
That relative extension is tame, with different exponent one.
Transitivity \eqref{eq:different-transitivity} gives
\[
 d_{\infty,n}
  =1+2\cdot\frac{3(3^n-1)}2
  =3^{n+1}-2.
\]
All places over infinity have the same data because the global
extension is Galois.
\end{proof}

We can now complete Theorem~\ref{thm:genus}. The ramification
indices at both zero and infinity exceed one, so these two points
are genuinely branched. Lemma~\ref{lem:zero} excludes any other
branch values. Put $N_n=|E_n|$. Over algebraically closed
constants there are $N_n/e_{0,n}$ places over zero and
$N_n/e_{\infty,n}$ places over infinity. The
Riemann--Hurwitz formula for the separable Galois cover therefore
reads
\begin{align*}
 2g_n-2
 &=N_n\left(-2+\frac{d_{0,n}}{e_{0,n}}
                 +\frac{d_{\infty,n}}{e_{\infty,n}}\right)\\
 &=N_n\left[-2+(1-2^{-n})+\left(\frac32-3^{-n}\right)\right]\\
 &=N_n\left(\frac12-2^{-n}-3^{-n}\right).
\end{align*}
This is \eqref{eq:genus}, and all assertions of
Theorem~\ref{thm:genus} follow.

For orientation, the first two instances of these proved formulas
are displayed in Table~\ref{tab:heights}. The table is a
specialization of the all-height proof, not evidence from which
the theorem is inferred.
\begin{table}[ht]
\centering
\caption{The first two geometric covers, obtained from the formulas.}
\label{tab:heights}
\begin{tabular}{crrrrrr}
\toprule
$n$ & $|E_n|$ & $e_{0,n}$ & $d_{0,n}$ &
$e_{\infty,n}$ & $d_{\infty,n}$ & $g_n$\\
\midrule
1 & 6 & 2 & 1 & 6 & 7 & 0\\
2 & 648 & 4 & 3 & 18 & 25 & 46\\
\bottomrule
\end{tabular}
\end{table}

The local/global distinction is already visible in the step from
height one to height two. The global quadratic and additive degrees
are $4$ and $27$, so the global degree is multiplied by $108$.
At infinity the quadratic step is split and the additive degree is
only $3$, so the local index is multiplied by $3$. The general
induction proves exactly this disparity with $3^n$ global additive
classes at height $n$ and one local class.
